\documentclass[11pt]{article}
\usepackage[margin=1in]{geometry}
\usepackage{amsmath,amssymb,amsthm}
\usepackage{booktabs,enumitem,array}
\usepackage[hidelinks]{hyperref}
\usepackage[expansion=false]{microtype}
\setlist{nosep,leftmargin=1.5em}

\newtheorem{theorem}{Theorem}
\theoremstyle{remark}
\newtheorem*{remark}{Remark}

\newcommand{\status}[1]{\hfill{\normalfont\small\textsc{[#1]}}}
\newcommand{\one}{\mathbf{1}}
\newcommand{\Wt}{\widetilde{W}}

\title{\textbf{v0.3f patch} \small(supersedes v0.3e)\\[2pt]
\large Final edits before merge: certificate semantics, the scope of the\\ symmetric statement, and the theorem header}
\author{18 September 2026}
\date{}

\begin{document}
\maketitle
\vspace{-2.5em}

All three edits are made. The centre $y$ is now explicitly not claimed to be a root; the symmetric statement is scoped to what holds for all $n$; and \texttt{PASS} now distinguishes the numerical cross-check from the rigorous certificate. The theorem header names the singleton-block model.

%======================================================================
\section*{1. The certificate is now standalone}

\texttt{certify\_standalone.py} contains no search, no random seed, no saved file, and no SciPy. It hardcodes the rational $\theta_1$, the 32-digit rational box centre $y$, and the radius $10^{-6}$, recomputes the exact target $p=F(\theta_1)$ with \texttt{Fraction} arithmetic, and checks the list below. The hard-coded centre $y$ is a 32-digit truncation of a high-precision approximation and is \emph{not} asserted to satisfy $F(y)=p$; the Krawczyk test certifies that a unique exact root $\theta_2\in X$ lies in the box around $y$, which is all the argument needs.
\begin{enumerate}[label=(C\arabic*)]
\item $\theta_1$ strictly inside the domain \hfill exact rational
\item $\det J(\theta_1)=-\tfrac{49846858308155727}{655360000000000000000000000000000000}\neq0$ \hfill exact rational
\item every point of $X=y\pm10^{-6}$ inside the domain \hfill exact rational
\item $\|\theta_1-y\|_\infty>10^{-6}$; the script prints the exact rational, $\approx0.0814612705830218$ \hfill exact rational
\item $K(X)\subset\operatorname{int}X$ and $\|I-YJ(X)\|_\infty<1$ \hfill interval
\end{enumerate}
No certificate predicate is evaluated in floating point: the comparisons are between exact rationals or between interval endpoints. The approximate inverse $Y$ is recomputed deterministically from $y$ and rationalized, so it carries no hidden dependence on the discovery run; only $\|I-YJ(X)\|_\infty<1$ matters, and any fixed $Y$ that achieves it would do.

(C5) runs twice, and the two runs have different status. The \texttt{mpmath.iv} pass routes some exact rationals through finite-precision \texttt{mpf} values before wrapping them as intervals, so it is a numerical cross-check; the Arb pass compares rigorous balls throughout and is the certificate. The script therefore prints
\begin{center}\small
\begin{tabular}{@{}ll@{}}
\toprule
\texttt{SANITY\_CHECK} & (C1)--(C4) and the \texttt{mpmath.iv} inclusion\\
\texttt{RIGOROUS\_CERTIFICATE} & (C1)--(C4), the Arb inclusion, and $\|I-YJ(X)\|_\infty<1$ in Arb\\
\bottomrule
\end{tabular}
\end{center}
and reports \texttt{RIGOROUS\_CERTIFICATE: NOT RUN} when \texttt{python-flint} is unavailable, never \texttt{PASS}. Both currently pass:
\begin{center}\small
\begin{tabular}{@{}lll@{}}
\toprule
Implementation & $K(X)\subset\operatorname{int}X$ & $\|I-YJ(X)\|_\infty$ bound\\
\midrule
\texttt{mpmath.iv} (40 dps), cross-check & true & $0.0516621376815982$\\
Arb balls via \texttt{python-flint} (200 bits), certificate & true & $0.0514087152550928\pm3.7\times10^{-61}$\\
\bottomrule
\end{tabular}
\end{center}
Arb's comparison operators are rigorous, so the inclusion test and $\|I-YJ(X)\|_\infty<1$ are decided inside ball arithmetic and both enter the \texttt{RIGOROUS\_CERTIFICATE} predicate. The two bounds differ in the third significant figure because the implementations round the enclosure differently; both are far below 1, and the inclusion margin is $5.2\times10^{-8}$ against a box radius of $10^{-6}$.

The discovery scripts (\texttt{witness.py}, \texttt{scan2.py}) stay in the appendix as provenance, clearly separated from the artifact. Nothing in the proof depends on \texttt{roots[0]} having been the second root.

%======================================================================
\section*{2. The symmetric statement}

\begin{theorem}[identification phase transition]
For the singleton-block relation-corrupted product model with $n$ conditionally independent reported views --- not the multi-verifier-per-check block-correlated model of Lemma~4(c):
\[
\boxed{\begin{array}{ll}
n\le3 &: \text{not locally identifiable, by dimension;}\\[1mm]
n=4 &: \text{generically locally identifiable, globally ambiguous on a nonempty open set;}\\[1mm]
n\ge5 &: \text{the aggregate model is generically globally identifiable;}\\[1mm]
\text{all }n &: (E^+,B^+)\ \text{and}\ (E^-,B^-)\ \text{remain internally non-identified.}
\end{array}}
\]
\end{theorem}

The fourth line is symmetric because in the exact model $E^+$ and $B^+$ both have $W^\star\equiv+\one$, while $E^-$ and $B^-$ both have $W^\star\equiv-\one$. Each pair is carried to a single product component by the nondifferential relation channel, so the internal splits $(\lambda_{E+},\lambda_{B+})$ and $(\lambda_{E-},\lambda_{B-})$ are never identified from the view law alone. Whenever the aggregate model \emph{is} identified --- in particular generically for $n\ge5$ --- the sums $S_+$ and $S_-$ are identified, but their internal splits are not. The blind-false-\emph{reject} pair is the mirror image of the blind-false-accept pair and costs nothing extra to state.

The weaker scoping is not a formality: at $n\le3$ nothing is identified, and at $n=4$ distinct global branches can disagree on the aggregate weights themselves. The certified pair of Section~2 is an instance --- $\theta_1$ has $S_+=3/10$ and $S_-=1/2$, while the root in the box has $S_+\approx0.2464$ and $S_-\approx0.4730$, gaps of $0.054$ and $0.027$, far outside the $10^{-6}$ box.

\begin{remark}
Because the paper is about false acceptance, say once after the theorem that subsequent sections work with $S_+$, and leave $S_-$ as the recorded symmetric case. A reader who wants the false-reject analysis then knows it exists and is identical.
\end{remark}

%======================================================================
\section*{3. Wording for the numerical scans}

Replace ``roughly half of random laws'' with a statement about the sampler:

\begin{quote}
In an exploratory scan, 10 of 20 targets generated from our parameter sampler yielded multiple feasible roots, and the search found up to three feasible roots for a single target.
\end{quote}

The scan draws $\theta\sim P_\Theta$ and sets $q=F(\theta)$, so the frequency is a property of $P_\Theta$ and of the multistart budget, not a volume on the space of laws; and multistart never proves that no further root exists. With the open-set theorem in hand these counts carry no load at all, so they belong in the appendix as intuition for why different targets behave differently.

Likewise the DSA result is a stronger numerical stress test on the corrected feasible region ($c_+=\lambda_0\rho$, $c_-=\lambda_0(1-\rho)$, $c_++c_-<1$), and 0/20 at $n=4$ and $n=5$ does not upgrade v0.1b's Proposition~2 beyond its current status: generic local identifiability, with global uniqueness still unproven. The wording in v0.1b already says exactly that and should not be touched.

Keep ``globally non-identifiable on a nonempty open set'' as the $n=4$ claim. ``Generically non-identifiable'' would require the discriminant and fiber-degree analysis and is not what the witness proves.

%======================================================================
\section*{4. The $n\ge5$ proof, in the order you suggest}

Write it entirely in the observed parameterization $(w,p_1,p_2,\eta)$ and keep $\alpha^{\pm}$ out until the end:
\begin{enumerate}
\item Group the views as $G_1=(W_1,W_2)$, $G_2=(W_3,W_4)$, $G_3=(W_5,\dots,W_n)$, with alphabets $4,4,2^{n-4}$. For $r=4$ classes the generic Kruskal ranks are $4,4,\min(4,2^{n-4})$, so $k_1+k_2+k_3\ge10=2r+2$ for every $n\ge5$, with equality at $n=5$. Kruskal's theorem then gives the four product components and their weights, up to permutation; this is the Allman--Matias--Rhodes route for mixtures of product distributions.
\item The tie-constrained submodel is not inside the exceptional set: at the rational witness $\det M_{12}=-63/8000$ and $\det M_{34}=-2783/116640$, both nonzero, and a fifth view with four pairwise-distinct component probabilities has Kruskal rank~2. A nonzero value at one point shows the relevant determinants are not identically zero on the submodel, so the condition holds generically there.
\item Resolve the permutation. Generically exactly one pair satisfies $q_{a,v}+q_{b,v}=1$ for all $v$; that is the atom pair. Orient it by $q_{A^+,v}>1/2>q_{A^-,v}$, and orient the remaining two by $p_{1v}>1/2>p_{2v}$. Any other pairing satisfying the complementarity relation lies in a proper algebraic subset.
\item Read off $\eta_v=q_{A^-,v}$, $p_{1v}$, $p_{2v}$ and the four weights, which include $S_+$ and $S_-$.
\item \emph{Corollary}, not part of the identification argument: $\alpha^+_v=\tfrac12+(p_{1v}-\tfrac12)/(1-2\eta_v)$ and $\alpha^-_v=\tfrac12+(\tfrac12-p_{2v})/(1-2\eta_v)$.
\end{enumerate}
Stating step~5 as a corollary keeps the theorem about the observed law and makes clear that the inversion needs $\eta_v\neq1/2$, which the domain already guarantees.

%======================================================================
\section*{5. What P1's conclusion now says}

Your one-sentence version is the right headline, and it is worth putting in the abstract almost verbatim: enough independent views eventually identify the entire latent measurement system, but no number of views identifies truth components that are observationally equivalent by construction. Redundancy identifies the apparatus; it cannot identify a distinction the measurement relation does not encode. That is the three-plate intuition stated as an identification theorem, and the gold-audit rule of Lemma~3 is what remains once it is accepted.

Deferring the generic fiber degree is right. It refines the middle line of the phase transition and explains the varying feasible-root counts, but the transition itself no longer depends on it.
\end{document}